% Part: computability
% Chapter: machines-computations
% Section: configuration

\documentclass[../../../include/open-logic-section]{subfiles}

\begin{document}

\olfileid{cmp}{tur}{con}
\olsection{Konfigurasi lan Komputasi}

\begin{explain}
Elinga nalika kita nelusuri konfigurasi mesin genep
ing bagean sadurunge. Mekanisme khayalan kang dumadi
saka pita, sirah maca-nulis, lan program mesin Turing
sejatine mung cara intuisi kanggo nggambarake
komputasi mesin Turing. Sacara formal, kita bisa
netepake komputasi mesin Turing kanthi input tartamtu
minangka rerangken \emph{konfigurasi}---lan konfigurasi
iku dhewe dumadi saka rerangken simbol (kang cocog
karo isi pita ing sawijining wektu komputasi),
wilangan kang nuduhake panggonan sirah maca-nulis,
lan sawijining kahanan. Kanthi iki, kita bisa
netepake apa kang diitung mesin Turing~$M$
nalika diwenehi input tartamtu.
\end{explain}

\begin{defn}[Konfigurasi]
\emph{Konfigurasi} mesin Turing $M = \tuple{Q, \Sigma, q_0,
\delta}$ yaiku tripel $\tuple{C, m, q}$ kanthi
\begin{enumerate}
\item $C \in \Sigma^*$ rerangken winates simbol saka $\Sigma$,
\item $m \in \Nat$ wilangan $< \len{C}$, lan
\item $q \in Q$
\end{enumerate}
Miturut intuisi, rerangken~$C$ yaiku isi pita
(simbol kabeh kothak saka kothak paling kiwa
nganti kothak pungkasan kang ora kosong utawa
wis tau ditekani), $m$~yaiku nomer kothak
kang lagi diwaca sirah maca-nulis (diwiwiti
saka $0$ minangka nomer kothak paling kiwa),
lan $q$ kahanan mesin saiki.
\end{defn}

\begin{explain}
Input kang bisa diwenehake marang mesin Turing
yaiku rerangken simbol, lumrahe rerangken kang
ngode sawijining wilangan kanthi cara tartamtu.
Konfigurasi wiwitan mesin Turing yaiku konfigurasi
nalika kita miwiti mesin Turing kanggo ngolah
input iku: pita ngemot panandha pucuk pita
kang langsung disusul input ing kothak-kothak
sisih tengene, sirah maca-nulis lagi maca
kothak paling kiwa saka input (yaiku kothak
ing sisih tengen panandha pucuk kiwa), lan
mekanisme ana ing kahanan wiwitan~$q_0$
kang wis ditemtokake.
\end{explain}

\begin{defn}[Konfigurasi wiwitan]
\emph{Konfigurasi wiwitan} $M$ kanggo input
$I \in \Sigma^*$ yaiku
\[
\begin{cases}
\tuple{\TMendtape \frown I, 1, q_0} & \text{yen input ora kosong},\\
\tuple{\TMendtape \frown \TMblank, 1, q_0} & \text{yen input kosong}.
\end{cases}
\]
Cathetan koreksi sumber OLPL-771 lan SOL6-F044: rumus saiki mbagi input kosong lan ora kosong. Ing kasus kosong, kothak blank kapisan ing sisih tengen panandha dilebokake ing rerangken konfigurasi supaya posisi sirah kang nomere siji sah. \end{defn}

\begin{explain}
% OLPL-146: Sumber beku nyebut input ing sisih kiwa panandha pucuk kiwa, nanging konfigurasi wiwitan lan andharan sadurunge nuduhake input ing sisih tengene.
Cathetan koreksi sumber OLPL-146: input ing sisih tengen panandha, dudu sisih kiwa kaya ukara sumber. Simbol~$\frown$ nuduhake \emph{pangandhengan
rerangken}---rerangken input diwiwiti langsung
ing sisih tengen panandha pucuk kiwa.
\end{explain}

\begin{defn}
Kita nyebut yen konfigurasi $\tuple{C, m, q}$
\emph{ngasilake konfigurasi $\tuple{C', m', q'}$
ing siji langkah} (miturut~$M$), yen lan mung yen
\begin{enumerate}
\item simbol kapangka-$m$ ing $C$ yaiku $\sigma$,
\item himpunan instruksi $M$ nemtokake $\delta(q, \sigma) =
  \tuple{q', \sigma', D}$,
\item simbol kapangka-$m$ ing $C'$ yaiku $\sigma'$, lan
\item
\begin{enumerate}
\item $D = L$ lan $m' = m - 1$ yen $m>0$, yen ora $m'=0$, utawa
\item $D = R$ lan $m' = m + 1$, utawa
\item $D = N$ lan $m' = m$,
\end{enumerate}
\item yen $m' = \len{C}$, mula $\len{C'} = \len{C} + 1$ lan simbol
  kapangka-$m'$ ing $C'$ yaiku~$\TMblank$. Yen ora, $\len{C'}=\len{C}$.
\item kanggo saben $i$ saengga $i < \len{C}$ lan $i \neq m$, $C'(i) = C(i)$,
\end{enumerate}
\end{defn}

\begin{defn}\ollabel{defn:run-output}
\emph{Lumakune $M$ kanthi input~$I$}
yaiku rerangken maksimal konfigurasi $C_i$ saka $M$,
kanthi $C_0$ konfigurasi wiwitan $M$
kanggo input~$I$, lan saben $C_i$ kang duwe konfigurasi sabanjure
ngasilake $C_{i+1}$ ing siji langkah. Cathetan panyunting SOL6-F047: rerangken iki bisa winates utawa tanpa wates. Yen winates, konfigurasi pungkasan ora duwe instruksi kang ditegesi; ora perlu penerus sawise mandheg.

Kita nyebut $M$ \emph{mandheg nalika nampa input $I$
sawise $k$ langkah} yen $C_k =
\tuple{C, m, q}$, simbol kapangka-$m$
ing~$C$ yaiku~$\sigma$, lan $\delta(q,
\sigma)$ ora ditegesi. Ing kahanan iku, yen panandha pucuk kiwa isih ana,
\emph{output}~$M$ kanggo input~$I$ yaiku~$O$,
ing ngendi $O$ rerangken simbol kang ora dipungkasi
dening~$\TMblank$ lan $C = \TMendtape \concat O \concat \TMblank^j$
kanggo sawijining~$j \in \Nat$. ($\TMblank^j$
yaiku rerangken $j$ simbol $\TMblank$.)
Cathetan cakupan SOL6-F048: definisi output iki mung menehi asil nalika panandha pucuk kiwa isih ana ing kothak nomer nol. Output bisa kosong; yen ora kosong, simbol pungkasan ora blank. Yen panandha kiwa wis kabusak, output ora ditegesi miturut konvensi iki, senajan mesin mandheg. \end{defn}

\begin{explain}
Miturut definisi iki, output~$O$ saka~$M$
mesthi dipungkasi simbol saliyane~$\TMblank$,
utawa, yen ing wektu~$k$ kabeh pita diisi
simbol~$\TMblank$ (kajaba simbol~$\TMendtape$
paling kiwa), $O$~minangka rerangken kosong.
\end{explain}

\end{document}
